% !TeX program = lualatex
\documentclass[10pt]{ltjsarticle} % titlepageオプションを外して警告を消去

% --- パッケージの読み込み ---
\usepackage[a4paper, vmargin=30mm, hmargin=25mm]{geometry} 
\usepackage{titlesec}      % セクション見出しのカスタマイズ
\usepackage{enumitem}      % 箇条書きのカスタマイズ
\usepackage{tabularx}      % 幅指定可能なテーブル
\usepackage{booktabs}      % 高品質な表の罫線
\usepackage{amssymb}       % 特殊記号（チェックボックス用など）
\usepackage{graphicx}      % 画像の挿入
\usepackage{url}           % URLの表示
\usepackage{pgfgantt}      % ガントチャートの作成
\usepackage{float}        % 図を強制配置するための [H]

% LuaLaTeXで日本語の文字化け・リンクの文字化けを防ぐ設定
\usepackage[luatex,unicode,pdfencoding=auto,psdextra]{hyperref} 

% --- レイアウト・装飾の設定 ---
\titleformat{\section}{\normalfont\Large\bfseries}{\thesection. }{0ex}{}
\titleformat{\subsection}{\normalfont\large\bfseries}{\thesubsection\ }{0ex}{}
\setlist[itemize]{leftmargin=2em, topsep=3pt, itemsep=3pt}

% --- ドキュメント開始 ---
\begin{document}

% テーブル内で長い文字列が重ならないようにする
\setlength{\extrarowheight}{2pt}

% --- タイトル・基本情報ヘッダー ---
\begin{center}
    {\huge\bfseries 作業ログ}
\end{center}
\vspace{10mm}

\noindent
\begin{tabularx}{\textwidth}{lX lX}
    \textbf{報告者:} & [近藤巧望] & \textbf{報告日:} & [6/30] \\
    \textbf{プロジェクト名:} & [長周期同期ストリーミング用いた通信負荷低減による低ジッター自動演奏システム] & & \\
\end{tabularx}

\vspace{5mm}
\hrule
\vspace{5mm}

% --- 2. 作業ログ ---
\section{作業ログ(詳細トラッキング)}
まず，担当箇所の計画前のガントチャートは図\ref{fig:pre-gc}，今週実施した作業の内容と状態は表\ref{tab:work-log}の通りである．

\begin{figure}[H]
    \centering
    \sloppy
    % 上: トロンボーン周り
    \begin{ganttchart}[
        hgrid,
        vgrid,
        x unit=0.45cm,
        y unit chart=0.7cm,
        time slot format=isodate,
        title height=1,
        bar/.style={fill=cyan!50, draw=cyan!80},
        bar height=0.6,
        group/.style={fill=gray!60, draw=gray!80},
        group height=0.2,
        group top shift=0.4
    ]{2026-05-20}{2026-06-17}

        \gantttitlecalendar{year, month=name, day} \\

        \ganttgroup{実装}{2026-05-21}{2026-06-17} \\
        \ganttbar{トロンボーンの音色}{2026-05-21}{2026-05-23} \\
        \ganttbar{トロンボーンの音の大きさ・長さ}{2026-05-24}{2026-05-26} \\
        \ganttbar{トロンボーンの音階作成}{2026-05-27}{2026-05-30} \\
        \ganttbar{トロンボーンのデバック}{2026-05-30}{2026-06-02} \\

    \end{ganttchart}

    \vspace{4mm}

    % 下: カスタネットとテスト
    \begin{ganttchart}[
        hgrid,
        vgrid,
        x unit=0.45cm,
        y unit chart=0.7cm,
        time slot format=isodate,
        title height=1,
        bar/.style={fill=cyan!50, draw=cyan!80},
        bar height=0.6,
        group/.style={fill=gray!60, draw=gray!80},
        group height=0.2,
        group top shift=0.4
    ]{2026-06-01}{2026-06-23}

        \gantttitlecalendar{year, month=name, day} \\

        \ganttbar{カスタネットの音色}{2026-06-05}{2026-06-07} \\
        \ganttbar{カスタネットの音の大きさ・長さ}{2026-06-08}{2026-06-10} \\
        \ganttbar{カスタネットの音階作成}{2026-06-11}{2026-06-13} \\
        \ganttbar{カスタネットのデバック}{2026-06-14}{2026-06-17} \\

        \ganttgroup{テスト}{2026-06-18}{2026-06-23} \\
        \ganttbar{各楽器の音テスト}{2026-06-18}{2026-06-18} \\
        \ganttbar{各楽器における演奏テスト}{2026-06-19}{2026-06-20} \\
        \ganttbar{最終演奏テスト}{2026-06-21}{2026-06-23} \\

    \end{ganttchart}
    \fussy
    \caption{計画前のGC}
    \label{fig:pre-gc}
\end{figure}

            

\begin{table}[h]
    \centering
    \caption{作業内容と状態}
    \small
    \begin{tabularx}{\textwidth}{|p{6em}|c|c|>{\raggedright\arraybackslash}X|c|>{\raggedright\arraybackslash}X|>{\raggedright\arraybackslash}X|}
    \hline
    \textbf{作業項目} & \textbf{開始予定日} & \textbf{完了日} & \textbf{状態} & \textbf{進捗率} & \textbf{依存関係・ブロッカー} & \textbf{コメント} \\ \hline
    自身が担当した実装作業で変更になった点について計画書を修正(3時間) & [6/24] & [6/26] & 完了 & [80\%] & [無] & [現状確認できた変更箇所を修正したため後に再度変更箇所を確認する．] \\ \hline
    トロンボーンの音確認テスト(3時間) & [6/26] & [6/28] & 完了 & [100\%] & [無] & [Arduinoから送信される周波数と実際に鳴る周波数の関係が正しいことを確認．] \\ \hline
    トロンボーンの正常発音数テスト(3時間) & [6/28] & [6/30] & 完了 & [100\%] & [無] & [シリアル受信した音数と譜面配列内の音数が一致することを確認．] \\ \hline
    \end{tabularx}
    \label{tab:work-log}
\end{table}

また，今週の作業内容を踏まえると計画後のガントチャートは図\ref{fig:after-gc}のようになると考えられる．
自身が担当した実装作業にて，クラス名や変数名などに細かな変更点があったため，計画書の修正を行った．
単体テストについてはまず，トロンボーンの音確認テストを行い，Arduinoから送信される周波数と実際に鳴る周波数の関係が正しいことを確認した．
Arduinoが送信する周波数のままではトロンボーンにしては音が高すぎることからprocessing内でオクターブを1下げる処理を行っているため，
双方の周波数が全く同じ数値にはならないが，音の高さの関係は正しいことを確認した．
同時に，音の波形表示を行った．
次に，トロンボーンの正常発音数テストを行い，シリアル受信した音数と譜面配列内の音数が一致することを確認し，
シリアル受信した情報をcsvファイルにまとめた．
\begin{figure}[H]
    \centering
    \sloppy
    % 上: トロンボーン周り
    \begin{ganttchart}[
        hgrid,
        vgrid,
        x unit=0.45cm,
        y unit chart=0.7cm,
        time slot format=isodate,
        title height=1,
        bar/.style={fill=cyan!50, draw=cyan!80},
        bar height=0.6,
        group/.style={fill=gray!60, draw=gray!80},
        group height=0.2,
        group top shift=0.4
    ]{2026-05-20}{2026-06-17}

        \gantttitlecalendar{year, month=name, day} \\

        \ganttgroup{実装}{2026-05-21}{2026-06-17} \\
        \ganttbar{トロンボーンの音色}{2026-05-21}{2026-05-23} \\
        \ganttbar{トロンボーンの音の大きさ・長さ}{2026-05-24}{2026-05-26} \\
        \ganttbar{トロンボーンの音階作成}{2026-05-26}{2026-05-26} \\
        \ganttbar{トロンボーンのデバック}{2026-05-27}{2026-06-05} \\

    \end{ganttchart}
    \fussy
    \caption{計画後のGC（前半）}
\end{figure}

\newpage

\begin{figure}[H]
    \centering
    \sloppy
    % 中盤: カスタネット・人形作成・計画書修正
    \begin{ganttchart}[
        hgrid,
        vgrid,
        x unit=0.45cm,
        y unit chart=0.7cm,
        time slot format=isodate,
        title height=1,
        bar/.style={fill=cyan!50, draw=cyan!80},
        bar height=0.6,
        group/.style={fill=gray!60, draw=gray!80},
        group height=0.2,
        group top shift=0.4
    ]{2026-06-05}{2026-06-26}

        \gantttitlecalendar{year, month=name, day} \\

        \ganttbar{カスタネットの音色}{2026-06-05}{2026-06-09} \\
        \ganttbar{カスタネットの音の大きさ・長さ}{2026-06-09}{2026-06-09} \\
        \ganttbar{カスタネットの音階作成}{2026-06-09}{2026-06-09} \\
        \ganttbar{カスタネットのデバック}{2026-06-10}{2026-06-17} \\
        \ganttbar{トロンボーン人形の作成}{2026-06-12}{2026-06-16} \\
        \ganttbar{カスタネット人形の作成}{2026-06-17}{2026-06-22} \\
        \ganttbar{計画書の修正}{2026-06-24}{2026-06-26} \\

    \end{ganttchart}
    \fussy
    \caption{計画後のGC（中盤：カスタネット・人形作成）}
\end{figure}

\begin{figure}[H]
    \centering
    \sloppy
    % テストフェーズ
    \begin{ganttchart}[
        hgrid,
        vgrid,
        x unit=0.45cm,
        y unit chart=0.7cm,
        time slot format=isodate,
        title height=1,
        bar/.style={fill=cyan!50, draw=cyan!80},
        bar height=0.6,
        group/.style={fill=gray!60, draw=gray!80},
        group height=0.2,
        group top shift=0.4
    ]{2026-06-18}{2026-07-08}

        \gantttitlecalendar{year, month=name, day} \\

        \ganttgroup{テスト}{2026-06-18}{2026-06-23} \\
        \ganttbar{トロンボーンの音確認テスト}{2026-06-26}{2026-06-28} \\
        \ganttbar{トロンボーンの正常発音数テスト}{2026-06-28}{2026-06-30} \\
        \ganttbar{各楽器における演奏テスト}{2026-07-01}{2026-07-01} \\
        \ganttbar{最終演奏テスト}{2026-07-08}{2026-07-08} \\

    \end{ganttchart}
    \fussy
    \caption{計画後のGC（テストフェーズ）}
    \label{fig:after-gc}
\end{figure}

\vspace{5mm}

% --- 3. 課題管理 ---
\section{課題管理}
\begin{table}[h]
    \centering
    \small
    \begin{tabularx}{\textwidth}{|c|>{\raggedright\arraybackslash}X|c|c|>{\raggedright\arraybackslash}X|c|c|}
    \hline
    \textbf{課題} & \textbf{発生原因} & \textbf{影響度} & \textbf{優先度} & \textbf{対策} & \textbf{期限} & \textbf{状態} \\ \hline
    特になし & [] &  &  & [] & [] &  \\ \hline
    \end{tabularx}
\end{table}
\noindent
\footnotesize{※影響度=プロジェクトへのインパクト \quad ※優先度=対応の緊急性}

\newpage

% --- 4. AI 利用状況と検証 ---
\section{AI 利用状況と検証(再現性重視)}

\subsection{利用内容}
\begin{itemize}
    \item \textbf{利用目的:} [実装]
    \item \textbf{入力(プロンプト要旨):}
    \begin{itemize}
        \item[$\circ$] [Arduinoのクライアントコードをテスト用に改良]
        \item[$\circ$] [Processing内での周波数計測・波形表示の実装]
        \item[$\circ$] [音確認テスト・正常発音数テストで得られた情報をまとめる．]
    \end{itemize}
    \item \textbf{出力概要:}
    \begin{itemize}
        \item[$\circ$] [示した要件通りに実装．正常に動作し，音が正しいか確認できた．]
        \item[$\circ$] [示した要件通りに実装．正常に動作し，送信した周波数と実際に鳴る周波数の関係が正しいことを確認できたとともに，音の波形表示も確認できた．]
        \item[$\circ$] [計測結果をわかりやすくまとめてもらった．]
    \end{itemize}
\end{itemize}

\subsection{検証方法}
\noindent{\footnotesize ※第三者が再現できるレベルで記述}
\begin{itemize}
    \item \textbf{検証手法:}
    \begin{itemize}
        \item[$\circ$] [テストコードが正しく動作するか実際に確認する．]
    \end{itemize}
    \item \textbf{参照資料:} \\
    \begin{itemize}
        \item[$\circ$] [Processing, Processing Serial library Reference, \url{https://processing.org/reference/libraries/serial/index.html}]
        \item[$\circ$] [compartmental.net, Minim, \url{https://code.compartmental.net/tools/minim/}]
    \end{itemize}
    \item \textbf{検証手順:}
    \begin{itemize}
        \item [$\circ$] 実際にコードを動作させ，音の確認や波形表示が正しく行われるか確認する．
        \item [$\circ$] 想定どおりの動作が確認できなければ，コードを確認．
        \item 不明な点があれば，公式ドキュメントを参照して確認する．
    \end{itemize}
\end{itemize}
\subsection{ファクトチェック結果}
\begin{itemize}
    \item \textbf{一致点:}
    \begin{itemize}
        \item[$\circ$] [送信した周波数より1オクターブ下の周波数でProcessing内で周波数計測・波形表示ができることを確認した．]
        \item[$\circ$] [シリアル受信した音数と譜面配列内の音数が一致することを確認した．]
    \end{itemize}
    \item \textbf{不一致・疑義:}
    \begin{itemize}
        \item[$\circ$] [特になし．]
    \end{itemize}
    \item \textbf{最終判断:} \quad $\square$ 妥当 \\ 
    \item \textbf{根拠:} テスト項目に必要な情報は一通り取得できたため．\\
\end{itemize}

\vspace{5mm}

% --- 5. 次回計画 ---
\section{次回計画}
\begin{itemize}
    \item \textbf{目標:}
    \begin{itemize}
        \item[$\circ$] [残りの単体テスト．]
        \item[$\circ$] [残りの結合テスト．]
    \end{itemize}
    \item \textbf{達成基準:}
    \begin{itemize}
        \item[$\circ$] [残りの単体テストが全て完了していること．]
        \item[$\circ$] [残りの結合テストが全て完了していること．]
    \end{itemize}
    \item \textbf{期限:}
    \begin{itemize}
        \item[$\circ$] [7/08]: 担当箇所含め全ての実装の完了
    \end{itemize}
\end{itemize}

\vspace{5mm}

% --- 6. その他 ---
\section{その他}
\begin{itemize}
    \item 連絡事項および特記事項は特になし．
\end{itemize}

\vspace{5mm}

% --- 7. 付録 ---
\section{付録}
\begin{itemize}
    \item \textbf{添付資料:}
    \begin{itemize}
        \item[$\circ$] HCK\_trombone.pde: [\url{https://github.com/sekiya-12/HCK_code/blob/main/HCK_trombone.pde}]
        \item[$\circ$] trombone\_tone\_test.ino: [\url{https://github.com/sekiya-12/HCK_code/blob/main/trombone_tone_test.ino}]
        \item[$\circ$] trombone\_serial\_test.ino: [\url{https://github.com/sekiya-12/HCK_code/blob/main/trombone_serial_test.ino}]
    \end{itemize}
    \item \textbf{添付範囲:} \\
    \indent 全体
\end{itemize}

\end{document}